\begin{titlepage}
\begin{abstract}
%
I code smell specifici del codice di machine learning possono degradare la qualità dei progetti, ma si sa poco su quanto la loro introduzione e la loro rimozione siano legate ad attività di refactoring. Questa tesi studia l'associazione tra smell e refactoring su 6350 occorrenze di smell tratte dal dataset NICHE e suddivise per dimensione del commit (small, medium, large). Per ciascuna occorrenza sono stati analizzati con PyRef il commit di introduzione e il commit di rimozione. Un refactoring è considerato valido se PyRef lo individua nello stesso file e nella stessa funzione dello smell; nel commit di rimozione, per deduzione, questo indica che il refactoring ha rimosso lo smell.
Soltanto una parte limitata delle occorrenze analizzabili ha un refactoring considerato valido ai fini dello studio e sommariamente, senza una divisione per dataset, possiamo affermare che è presente un  refactoring valido nel 11\%–18\% dei commit di rimozione e nell’9\%–12\% di quelli di introduzione. L’analisi della correlazione di Spearman tra la presenza di un refactoring valido e il numero di smell presenti nel commit, calcolata separatamente per introduzione e rimozione su tutti e tre i subset, mostra un’associazione negativa debole-moderata, con ρ compreso tra −0,10 e −0,23 per l’introduzione e tra −0,11 e −0,33 per la rimozione. L’associazione è più marcata per la rimozione, dove risulta statisticamente significativa nei subset large (ρ = −0,25) e medium (ρ = −0,33), mentre nel subset small non lo è (ρ = −0,11, p = 0,20). Questo indica che i refactoring sono associati, seppur in maniera relativa, all’introduzione e/o alla rimozione di smell, come mostrato dai valori di ρ e dai p-value sopra riportati
%
\\[1cm]
\end{abstract}
\end{titlepage}